\documentclass[11pt, a4paper]{article}

\usepackage[a4paper, top=2.5cm, bottom=2.5cm, left=2cm, right=2cm]{geometry}
\usepackage{amsmath}
\usepackage{amssymb}
\usepackage[colorlinks=true, linkcolor=blue, urlcolor=blue, citecolor=blue]{hyperref}

\title{Theory of Everything - Volume 10 \\ \Large The Standard Model of Particle Physics}
\author{Omega Protocol}
\date{\today}

\begin{document}

\maketitle

\section*{Layman's Summary}
The Standard Model is the recipe book for the universe, describing all the basic ingredients (quarks, electrons) and the forces that stick them together. Volume 10 explains *why* this specific recipe exists. Instead of just accepting the particles we observe, we show how they naturally emerge from the fundamental geometry and information flow of space itself.

\section{Gauge Symmetry Emergence (Axiom 13)}
Forces like electromagnetism and the strong nuclear force are based on mathematical rules called "gauge symmetries" (specifically SU(3), SU(2), and U(1)). We prove that these aren't arbitrary rules handed down from nowhere. Instead, they naturally bubble up from the way the underlying quantum information network twists and vibrates (the "modular Dirac operator").

\section{The Higgs Mechanism (Theorem)}
Why do particles have mass? The famous Higgs boson gives particles mass by creating a field that "drags" on them. In our theory, we show that this isn't an accident. As the early universe cooled down, the underlying informational network had to settle into a stable state. This "cooling" process mathematically forces the creation of the Higgs mechanism, spontaneously giving particles their mass.

\section{Fermion Generations (Theorem)}
There are exactly three "families" or generations of matter particles (like the electron, muon, and tau). Physicists have long wondered why it's three and not two or four. We prove that this number is determined by the specific shape and topology of the emergent spacetime manifold. The geometry of the universe dictates the families of matter that can exist within it.

\end{document}
